The hospital findings concern one synthetic facility and declared parameter
ranges. Binary service availability, a 72-hour horizon, and assumed control
effects limit transport to clinical capacity, patient harm, regional spread,
or months of recovery. The external-validation registry contains nine
benchmarks: one numeric failure, one loose numeric consistency result, two
qualitative failures, three not scored, and two structurally unaddressable.
These comparisons do not validate procurement recommendations.

The interval certificate is exact for the declared Cartesian box, conditional
on five fixed estimated objectives in the hospital application. It does not
account for their Monte Carlo error or for structural uncertainty. Correlated
specification sets can make its bounds conservative. The structural analysis
uses selected candidates, and only one structural choice is varied. The
bootstrap conditions on the observed bank and frozen shortlist; ordinary
percentile coverage for discontinuous Pareto selection is not established.
The family-adjusted mean intervals are approximate.

The engineering problems have independent published definitions but were
evaluated in the same workflow. Their finite sampled candidates and synthetic
bounds test transfer and implementation, not clinical realism or independently
conducted replication. Historical assertion cases were selected after review,
so they cannot establish unseen-error detection accuracy. An omitted or weak
assertion can leave an incorrect sentence unchecked.

The study uses synthetic trials and public aggregate sources, with no patient
records, participant recruitment, or live-system experiments. The artifact
preserves failed comparisons and historical source-state limitations so that
readers can distinguish a verified computational relationship from an
empirically identified mechanism.
